\documentclass[11pt, letterpaper]{article}

% --- PAQUETES ---
\usepackage[utf8]{inputenc}
\usepackage[T1]{fontenc}
\usepackage[spanish, es-noshorthands]{babel}
\usepackage[top=2cm, bottom=2.2cm, left=2.2cm, right=2.2cm]{geometry}
\usepackage{graphicx}
\usepackage{amsmath}
\usepackage{amssymb}
\usepackage{enumitem}
\usepackage[table]{xcolor}
\usepackage{fancyhdr}
\usepackage{eso-pic}
\usepackage{multicol}
\usepackage{array}
\usepackage{pifont}
\usepackage[most]{tcolorbox}
\usepackage{tikz}

% --- FUENTE SISTEMA ---
\usepackage{helvet}
\renewcommand{\familydefault}{\sfdefault}

% Fondo institucional Liceo V.A.L.
\IfFileExists{../../FondoVAL.pdf}{
	\AddToShipoutPictureBG{
		\AtPageLowerLeft{
			\includegraphics[width=\paperwidth,height=\paperheight,keepaspectratio=false]{../../FondoVAL.pdf}
		}
	}
}{}

% --- CONFIGURACIÓN DE PÁGINA ---
\pagestyle{fancy}
\fancyhf{}
\fancyhead[L]{\footnotesize \color{gray} LICEO V.A.L. -- DEPARTAMENTO DE MATEMÁTICAS}
\fancyhead[R]{\footnotesize \color{gray} MATEMÁTICAS 6º -- AC}
\fancyfoot[C]{\footnotesize \color{gray} Página \thepage}
\renewcommand{\headrulewidth}{0.4pt}
\renewcommand{\footrulewidth}{0pt}

% --- INTERRUPTOR DE RESPUESTAS ---
\newif\ifshowanswers
\showanswersfalse % Cambiar a \showanswerstrue para el solucionario

% ===================================================================
% MACROS GRÁFICAS PARA ADAPTACIÓN CURRICULAR (AC) - MODO AHORRO TINTA
% ===================================================================

% CORRECCIÓN: Se define con [2][3cm] para aceptar un argumento opcional y uno obligatorio
\newcommand{\boxfill}[2][3cm]{%
	\ifshowanswers
	\tcbox[on line,enhanced,size=minimal, colframe=black, colback=black!10, boxrule=1pt, arc=2mm, boxsep=4pt, left=4pt, right=4pt, top=4pt, bottom=4pt]%
	{\makebox[#1]{\vphantom{Hg}\bfseries\color{black}#2}}%
	\else
	\tcbox[on line,enhanced,size=minimal, colframe=black!50, colback=white, boxrule=1pt, arc=2mm, boxsep=4pt, left=4pt, right=4pt, top=4pt, bottom=4pt]%
	{\makebox[#1]{\vphantom{Hg}}}%
	\fi
}

% CORRECCIÓN: Se cambia de [1] a [2]
\newcommand{\billete}[2]{%
	\begin{tcolorbox}[enhanced, colback=white, colframe=black!40, arc=3mm, boxrule=0.8pt, left=4pt, right=4pt, top=4pt, bottom=4pt]
		\begin{minipage}[c]{0.48\linewidth}
			\centering
			\IfFileExists{billete_#1_frente.png}{%
				\includegraphics[width=\linewidth,height=1.8cm, keepaspectratio=true]{billete_#1_frente.png}\\[2pt]%
			}{%
				\tcbox[colback=white,colframe=black!40,arc=2mm]{\small Billete (Frente)}\\[2pt]%
			}%
			{\tiny\textcolor{black!70}{\textsc{Frente}}}
		\end{minipage}\hfill
		\begin{minipage}[c]{0.48\linewidth}
			\centering
			\IfFileExists{billete_#1_reverso.png}{%
				\includegraphics[width=\linewidth,height=1.8cm, keepaspectratio=true]{billete_#1_reverso.png}\\[2pt]%
			}{%
				\tcbox[colback=white,colframe=black!40,arc=2mm]{\small Billete (Reverso)}\\[2pt]%
			}%
			{\tiny\textcolor{black!70}{\textsc{Reverso}}}
		\end{minipage}
		\vspace{2pt}
		\begin{center}
			{\small\bfseries\textcolor{black}{#2}}
		\end{center}
	\end{tcolorbox}%
}

% CORRECCIÓN: Se cambia de [1] a [2]
\newcommand{\moneda}[2]{%
	\begin{tcolorbox}[enhanced, colback=white, colframe=black!40, arc=3mm, boxrule=0.8pt, left=4pt, right=4pt, top=4pt, bottom=4pt]
		\begin{minipage}[c]{0.48\linewidth}
			\centering
			\IfFileExists{moneda_#1_cara.png}{%
				\includegraphics[width=1.5cm,height=1.5cm, keepaspectratio=true]{moneda_#1_cara.png}\\[2pt]%
			}{%
				\tcbox[colback=white,colframe=black!40,arc=2mm]{\small Moneda (Cara)}\\[2pt]%
			}%
			{\tiny\textcolor{black!70}{\textsc{Cara}}}
		\end{minipage}\hfill
		\begin{minipage}[c]{0.48\linewidth}
			\centering
			\IfFileExists{moneda_#1_sello.png}{%
				\includegraphics[width=1.5cm,height=1.5cm, keepaspectratio=true]{moneda_#1_sello.png}\\[2pt]%
			}{%
				\tcbox[colback=white,colframe=black!40,arc=2mm]{\small Moneda (Sello)}\\[2pt]%
			}%
			{\tiny\textcolor{black!70}{\textsc{Sello}}}
		\end{minipage}
		\vspace{2pt}
		\begin{center}
			{\small\bfseries\textcolor{black}{#2}}
		\end{center}
	\end{tcolorbox}%
}

\newtcolorbox{calculadora}{
	enhanced, colback=white, % Fondo blanco para no gastar tinta
	colframe=black, arc=4mm, boxrule=1.2pt, left=8pt, right=8pt, top=6pt, bottom=6pt, before skip=4pt, after skip=4pt
}


% --- MACROS PARA CASILLAS DE VERIFICACIÓN Y BOTONES ---
\newcommand{\casillaVacia}{%
	\tcbox[on line, enhanced, size=minimal, colframe=black!60, colback=white, boxrule=1pt, arc=1mm, left=2pt, right=2pt, top=2pt, bottom=2pt]{\makebox[0.4cm]{\vphantom{X}}}%
}
\newcommand{\casillaX}{%
	\tcbox[on line, enhanced, size=minimal, colframe=black, colback=black!15, boxrule=1.5pt, arc=1mm, left=2pt, right=2pt, top=2pt, bottom=2pt]{\makebox[0.4cm]{\vphantom{X}\textbf{\textsf{X}}}}%
}
\newcommand{\btn}[1]{%
	\tcbox[on line, enhanced, size=minimal, colframe=black!50, colback=white, boxrule=1pt, arc=2mm, left=8pt, right=8pt, top=4pt, bottom=4pt]{\textbf{#1}}%
}
\newcommand{\btnX}[1]{%
	\tcbox[on line, enhanced, size=minimal, colframe=black, colback=black!15, boxrule=1.5pt, arc=2mm, left=8pt, right=8pt, top=4pt, bottom=4pt]{\textbf{#1}}%
}
\begin{document}
	
	\begin{center}
		{\Large \textbf{LICEO V.A.L. (Vida, Amor, Luz)}}\\[0.1cm]
		{\large \textbf{DEPARTAMENTO DE MATEMÁTICAS -- GRADO 6º}}\\[0.1cm]
		{\large \textbf{Unidad 11: Manejo de Dinero (Pago, Vueltas y Facturas)}}\\[0.1cm]
		{\small \hfill \textbf{Código: 06M--12 AC}}
	\end{center}
	
	\vspace{0.2cm}
	\noindent \textbf{Nombre del Estudiante:} \underline{\hspace{8.5cm}} \textbf{Fecha:} \underline{\hspace{3cm}}
	\vspace{0.3cm}
	
	\begin{enumerate}[leftmargin=*, itemsep=0.8cm]
		
		% PUNTO 1
		\item \textbf{(25\%) Para comprar cada producto necesitas entregar un billete de valor igual o mayor a su precio. Marca con una X el billete adecuado con el que puedes pagar cada compra:}
		
		\begin{center}
			\renewcommand{\arraystretch}{1.8}
			\arrayrulecolor{black!60}
			\setlength{\arrayrulewidth}{1pt}
			
			\begin{tabular}{| >{\centering\arraybackslash}m{3.5cm} | >{\centering\arraybackslash}m{2.5cm} | m{8cm} |}
				\hline
				\rowcolor{black!15}
				\textbf{Producto a Comprar} & \textbf{Precio} & \multicolumn{1}{c|}{\textbf{¿Con qué billete debes pagar?}} \\
				\hline
				Una pizza grande & \textbf{\large \$18.000} & 
				\ifshowanswers 
				\casillaVacia\ Billete de \$10.000 \quad \casillaX\ Billete de \$20.000
				\else 
				\casillaVacia\ Billete de \$10.000 \quad \casillaVacia\ Billete de \$20.000 
				\fi \\
				\hline
				\rowcolor{black!04}
				Un balón de fútbol & \textbf{\large \$35.000} & 
				\ifshowanswers 
				\casillaVacia\ Billete de \$20.000 \quad \casillaX\ Billete de \$50.000
				\else 
				\casillaVacia\ Billete de \$20.000 \quad \casillaVacia\ Billete de \$50.000 
				\fi \\
				\hline
				Un helado artesanal & \textbf{\large \$3.500} & 
				\ifshowanswers 
				\casillaVacia\ Billete de \$2.000 \quad \casillaX\ Billete de \$5.000
				\else 
				\casillaVacia\ Billete de \$2.000 \quad \casillaVacia\ Billete de \$5.000 
				\fi \\
				\hline
				\rowcolor{black!04}
				Unos tenis deportivos & \textbf{\large \$85.000} & 
				\ifshowanswers 
				\casillaVacia\ Billete de \$50.000 \quad \casillaX\ Billete de \$100.000
				\else 
				\casillaVacia\ Billete de \$50.000 \quad \casillaVacia\ Billete de \$100.000 
				\fi \\
				\hline
			\end{tabular}
		\end{center}
		
		% PUNTO 2
		\item \textbf{(25\%) Compara el dinero que tienes en la mano con la cuenta a pagar. Marca la opción correcta (SÍ o NO):}
		
		\begin{center}
			\renewcommand{\arraystretch}{2} % Aumentado para que los botones quepan bien
			\arrayrulecolor{black!60}
			\setlength{\arrayrulewidth}{1pt}
			
			\begin{tabular}{| >{\centering\arraybackslash}m{3.5cm} | >{\centering\arraybackslash}m{3cm} | >{\centering\arraybackslash}m{4.5cm} | >{\centering\arraybackslash}c |}
				\hline
				\rowcolor{black!15}
				\textbf{Dinero que llevo} & \textbf{Cuenta a pagar} & \textbf{Comparación} & \textbf{¿Me alcanza?} \\
				\hline
				\textbf{\$10.000} en billete & \textbf{\$8.500} & \$10.000 es \textbf{MAYOR} que \$8.500 & 
				\ifshowanswers \btnX{SÍ} \quad \btn{NO} \else \btn{SÍ} \quad \btn{NO} \fi \\
				\hline
				\rowcolor{black!04}
				\textbf{\$5.000} en billete & \textbf{\$7.200} & \$5.000 es \textbf{MENOR} que \$7.200 & 
				\ifshowanswers \btn{SÍ} \quad \btnX{NO} \else \btn{SÍ} \quad \btn{NO} \fi \\
				\hline
				\textbf{\$20.000} en billete & \textbf{\$15.000} & \$20.000 es \textbf{MAYOR} que \$15.000 & 
				\ifshowanswers \btnX{SÍ} \quad \btn{NO} \else \btn{SÍ} \quad \btn{NO} \fi \\
				\hline
				\rowcolor{black!04}
				\textbf{\$2.000} en billete & \textbf{\$3.800} & \$2.000 es \textbf{MENOR} que \$3.800 & 
				\ifshowanswers \btn{SÍ} \quad \btnX{NO} \else \btn{SÍ} \quad \btn{NO} \fi \\
				\hline
			\end{tabular}
		\end{center}
		\newpage
		% PUNTO 3
		\item \textbf{(25\%) Calcula las vueltas que debe entregarte el cajero realizando la resta en tu calculadora:}
		
		\begin{enumerate}[label=\alph*), itemsep=0.4cm]
			\item Pagas una cuenta de \textbf{\$14.500} con un billete de \textbf{\$20.000}:
			\begin{center}
				\begin{calculadora}
					\textcolor{black}{\small En la calculadora:}\\[2pt]
					\textcolor{black}{\large\bfseries 20000 \;-\; 14500 \;=}
				\end{calculadora}
				\vspace{3pt}
				{\large \textbf{MIS VUELTAS = \$}\; \boxfill[3cm]{\$5.500} }
			\end{center}
			
			\item Pagas una cuenta de \textbf{\$32.000} con un billete de \textbf{\$50.000}:
			\begin{center}
				\begin{calculadora}
					\textcolor{black}{\small En la calculadora:}\\[2pt]
					\textcolor{black}{\large\bfseries 50000 \;-\; 32000 \;=}
				\end{calculadora}
				\vspace{3pt}
				{\large \textbf{MIS VUELTAS = \$}\; \boxfill[3cm]{\$18.000} }
			\end{center}
		\end{enumerate}
		
		% PUNTO 4
		\item \textbf{(25\%) Observa el tiquete de compra de la droguería y responde las preguntas leyendo los datos directamente de la factura:}
		
		\begin{center}
			\begin{minipage}[c]{0.42\linewidth}
				\begin{tcolorbox}[enhanced, colback=white, colframe=black, arc=2mm, boxrule=1.2pt]
					\centering
					{\bfseries DROGUERÍA LICEO V.A.L.}\\[2pt]
					{\tiny NIT: 900.123.456-1 \quad Tel: 300 123 4567}\\[4pt]
					\hrule\vspace{4pt}
					\small
					\begin{tabular}{lr}
						1x Alcohol Antiséptico & \$4.500 \\
						1x Tapabocas x3        & \$2.000 \\
					\end{tabular}\\[4pt]
					\hrule\vspace{4pt}
					\begin{tabular}{lr}
						\textbf{TOTAL A PAGAR:} & \textbf{\$6.500} \\
						EFECTIVO RECIBIDO:      & \$10.000 \\
						\textbf{VUELTAS:}       & \textbf{\$3.500} \\
					\end{tabular}\\[4pt]
					\hrule\vspace{4pt}
					{\tiny ¡Gracias por su compra!}
				\end{tcolorbox}
			\end{minipage}\hfill
			\begin{minipage}[c]{0.54\linewidth}
				\begin{enumerate}[label=\alph*), itemsep=0.4cm]
					\item ¿Cuál es el \textbf{TOTAL A PAGAR} de la compra? \\
					\textbf{Respuesta: \$}\; \boxfill[2.5cm]{\$6.500}
					
					\item ¿Cuánto \textbf{EFECTIVO} recibió el cajero? \\
					\textbf{Respuesta: \$}\; \boxfill[2.5cm]{\$10.000}
					
					\item ¿De cuánto fueron las \textbf{VUELTAS}? \\
					\textbf{Respuesta: \$}\; \boxfill[2.5cm]{\$3.500}
				\end{enumerate}
			\end{minipage}
		\end{center}
		
	\end{enumerate}
	
	\vfill
	\begin{center}
		\textit{\color{black!70} ``Revisar mi factura y contar mis vueltas me ayuda a cuidar mi dinero.'' --- Liceo V.A.L.}
	\end{center}
	
\end{document}